No one set out to write it. It began as an ordinary applied project with an ordinary ambition:
show that anchoring a neural quantile model to a classical VaR prior beats the unanchored model.
The intended paper was the one that gets written thousands of times a year --- \emph{our method wins,
here is the table.}

The starting state is worth describing precisely, because it is unremarkable and that is the
point:

\begin{itemize}
\tightlist
\item
  The method was called a \textbf{physics-informed neural network}. It was not. No monotonicity or
  sub-additivity constraint was ever implemented; the mechanism was an L2 penalty toward a
  rolling parametric prior. The name came from the family of ideas the work was inspired by,
  and it survived because nothing in the workflow required anyone to check it against the code.
  \textbf{The paper's own title was wrong before any result was.}
\item
  Hyper-parameters were chosen by running the pipeline and looking at the output. There was no
  validation block, so they were selected on the test set.
\item
  Each configuration ran on \textbf{one seed}.
\item
  Models were ranked on breach counts and ``capital reserved'' rather than on a strictly
  consistent loss.
\item
  The GARCH benchmark carried a standardized-\emph{t} scaling error that inflated its VaR by
  \(\nu/(\nu-2)\) --- that is, \textbf{the baseline we were beating was handicapped} (\Cref{sec:defects}).
\end{itemize}

Every item on that list is a decision someone made while trying to do good work, and none would
have been caught by peer review of the resulting manuscript, because none is visible in a
manuscript. Fixing them was not a matter of discovering misconduct. It was a matter of a
protocol asking questions the original workflow had no place to put.

What followed is the substance of this paper: each correction moved the result toward the null,
and after four of them there was no result left. The four withdrawals below are not a catalogue
of things other researchers do wrong. \textbf{They are what one careful project looked like when it was
finally instrumented well enough to see itself.}
